\section{Related work}
\label{sec:related}

\paragraph{How much a loop is worth.}
Sharing parameters across layers is an old way to trade compute for parameters
\citep{dehghani2019universal,lan2019albert,bae2025relaxed}. \citet[App.~D.2, Fig.~17]{kaplan2020scaling} found that transformers
which re-use parameters ``perform slightly better when comparing models with equal parameter count, but slightly worse
when accounting for reuse and comparing per FLOP''; our results agree in both directions and add how the balance moves
with model size and data. The effective-parameter form we test, $\Neff = \Nonce + r^{\varphi}\Nrec$, and our choice
of two unshared layers before and after the loop come from \citet{schwethelm2026how}, who estimate $\varphi = 0.46$
(95\% interval 0.41--0.53) on FineWeb-Edu with up to eight loops and find that backpropagating through only the last
$\lceil r/2 \rceil$ loops, as \citet{prairie2026parcae} do, lowers it from 0.45 (their baseline refitted on the same budgets) to
0.38, in a comparison that spends the
saved FLOPs on more tokens; our pre-registration was built to be comparable with theirs, and its H1 tests that result
with the two accountings separated. The prelude--core--coda layout with an input injection that maps the concatenated
state and prelude output is that of Huginn \citep{geiping2025scaling}. Parcae \citep{prairie2026parcae} finds that looping lowers perplexity at matched parameters and tokens, and that
at a fixed FLOP budget looping and data should be increased in tandem; \citet{lee2026sparse} that dense looped models
scale worse than non-looped ones at equal FLOPs; \citet{chen2026growth} that tying the weights of a deeper model costs a
roughly constant factor of compute in single-epoch training, while with repeated data the best loop count grows with
compute; for dense looped models and looped mixtures of experts alike, a bounded form extrapolates better than
$r^{\varphi}$ to 16 loops \citep{chen2026moe}, and in single-seed ablations of a compute-matched looped mixture of experts, looping the middle
half of the layers beat looping all of them by less than 0.01 nats \citep{wang2026smelt}. In a controlled appendix of
Ouro \citep{zhu2025scaling} a dense model of the same depth beats the looped one in every comparison, and plain
recursion is worse than a vanilla model in validation loss at matched compute in \citet{bae2025mixture}. A 1.4B looped model trained on
310B tokens (210B of pretraining, about 13 passes over a 16B-token corpus, then 100B of mid-training) beats a dense
model with the same parameters and data on all 12 of its benchmarks at about four times the training compute per
token, while a 3.9B dense model that matches its inference compute at three loops is ahead or level on every
benchmark \citep{marchenko2026closing}. When data is limited, loops may pay most: recursive transformers beat standard ones on downstream evaluations
when pretraining on 10M--100M words \citep{gulbahar2026squeezing}. None of these studies varies model size and tokens per
parameter separately with several seeds, which is what our grid adds; Section~\ref{sec:multiplier} restates their
equal-parameter results in a common unit.

\paragraph{What the loop computes.}
Huginn \citep{geiping2025scaling}, trained with a random loop count, has latent trajectories that, in qualitative
examples, settle into fixed points or orbits or drift steadily in one direction; \citet{logan2026pertoken} reports per-token convergence to a fixed point in a small depth-recurrent
model; deep equilibrium models make the fixed point the definition of the output \citep{bai2019deep}, and those whose
fixed point does not depend on the starting state make better use of extra test-time computation \citep{anil2022path}.
\citet{saunshi2025reasoning} find that the gains of looped models concentrate on reasoning tasks. We read the state after
every loop with the model's own output layers and with learned maps in front of them, as the logit lens
\citep{nostalgebraist2020logit} and the tuned lens \citep{belrose2023tuned} do for the layers of a dense model, for three
training regimes.

\paragraph{Serial reasoning and extrapolation.}
Recurrent networks solve harder instances than they were trained on by running more iterations, and can
``overthink'' \citep{schwarzschild2021can,bansal2022endtoend}. Looped transformers with an input-dependent number of
loops generalise to longer inputs on algorithmic tasks \citep{fan2025looped}, and on top of position embeddings made for
the task, recurrent layers with input injection help arithmetic generalise further to longer numbers
\citep{mcleish2024arithmetic}. Recurrent-depth transformers extrapolate to
deeper chains of composed facts with more loops at test time (trained on up to 12 hops, they reach 14--19), while
vanilla transformers of equal parameters or equal FLOPs do not \citep{kohli2026loop}; and the positions a loop resolves
per iteration scale with the training length divided by the training loop count, within a task-dependent range
\citep{zhang2026recurrence}. Shallow
transformers can also learn shortcuts for serial computations whose steps compose associatively
\citep{liu2022shortcuts}. Our synthetic results agree with these studies; we add probes of every
intermediate value after every loop, the effect of the learning rate on whether the task is learned within the training budget, a
controlled comparison of stopping rules, and a comparison, on the same problems, with a chain of thought
\citep{wei2022chain}, which adds serial computation to a fixed-depth transformer \citep{merrill2024expressive} and is
learned here from scratch on the written intermediate values, as scratchpad training does for pre-trained models
\citep{nye2021show}, and with pause tokens \citep{goyal2024think}, which give it extra parallel computation; filler tokens of this kind help on some
algorithmic tasks \citep{pfau2024dot}. RELAY \citep{yu2025relay} aligns loop iterations with chain-of-thought steps and
uses a looped model's length generalisation to produce chains for longer problems, and continuous thoughts
\citep{hao2025training} feed the hidden state back as input. Transformers that write tokens are known to generalise
poorly to inputs longer than those seen in training \citep{zhou2024transformers,cho2024position}, which is how our
chain-of-thought baseline fails.

\paragraph{When to stop.}
Adaptive computation time \citep{graves2016adaptive}, used by Universal Transformers, and PonderNet
\citep{banino2021pondernet} learn a halting distribution during training; Mixture-of-Recursions learns a recursion depth
per token \citep{bae2025mixture}. \citet{fan2025looped} report results with the ground-truth number of steps and propose,
as an alternative, the loop count up to a maximum at which the model is most confident in its own output, chosen per test
set or per problem; stopping when a confidence threshold is first met is mentioned there as a variant. We compare answer-free stopping rules applied after
training (thresholds on the state change, on the prediction and on the confidence, and votes over loop counts) on the
same checkpoints across training regimes, with the rules and thresholds fixed in advance and selected on a separate
problem set.

\paragraph{Cost at inference.}
Without sharing, a looped model's KV cache and decoding time grow with the loop count: two loops take about twice the
KV memory and 1.96 times the latency in \citet{wu2025parallel}, and plain looped decoding steps take 2.3, 3.4 and 5.7
times as long at 2, 3 and 5 loops in \citet{monea2026shared}. Ouro \citep{zhu2025scaling} shares the cache across loops during
decoding only, since processing the prompt needs a cache per loop; \citet{huang2026towards} share it and find that full
backpropagation through a fixed depth prevents sharing; and \citet{monea2026shared} pretrain looped models in which
only the first loop writes a cache and later loops read it while keeping a short window of their own, and find that
this costs no quality. We test sharing
only as a change made after training (Section~\ref{sec:cost}).
